\chapter{Root Cause Analysis Engine}
\label{ch:root_cause_analysis}

\section{Motivation}
\label{sec:rca_motivation}

Detecting anomalies in distributed systems is necessary but insufficient for 
maintaining reliability. In microservice-based architectures, failures often 
propagate across service boundaries, generating secondary anomalies that obscure 
the original fault \cite{munoz2017impact, zhang2019tracing}. As a result, simply 
flagging abnormal behavior does not provide operators with enough information to 
act effectively.

Root cause analysis (RCA) addresses this limitation by identifying the service or 
component most likely responsible for the observed anomaly pattern. Effective RCA 
reduces Mean Time to Recovery (MTTR), supports Service Level Agreement (SLA) 
compliance, and improves operational efficiency by prioritizing the true source 
of failure \cite{fu2020root}. Let $\mathcal{S} = \{S_1, \dots, S_n\}$ denote the 
set of services and $\mathcal{A} \subseteq \mathcal{S}$ the subset of anomalous 
services at time $t$. The RCA problem can be formalized as:

\[
R: \mathcal{A} \mapsto \mathcal{S}_{root} \subseteq \mathcal{S},
\]

where $\mathcal{S}_{root}$ represents the most likely root cause candidates.

\section{Distributed Tracing}
\label{sec:rca_tracing}

\subsection{Trace Collection}
\label{subsec:rca_trace_collection}

Distributed tracing provides an observational basis for RCA by capturing 
inter-service request flows. A single request is associated with a trace 
identifier and decomposed into spans representing operations executed across 
services \cite{sigelman2010dapper}. Modern tracing ecosystems such as 
OpenTelemetry, Jaeger, and Zipkin enable this form of end-to-end request 
visibility \cite{opentelemetry2023, jaeger2024, zipkin2024}.

A trace can be represented as:

\[
\tau = \{sp_1, sp_2, \dots, sp_m\},
\]

where each span corresponds to a unit of work executed within a service. A set of 
traces $\mathcal{T}$ allows the system to infer service dependencies, latency 
patterns, and anomaly propagation paths.

\subsection{Trace Structure: Spans and Context}
\label{subsec:rca_trace_structure}

A span describes an operation within a service and typically contains a span 
identifier, parent span identifier, service name, operation name, timestamps, 
status, and optional metadata. Parent--child span relationships define the 
execution structure of a distributed request:

\[
Parent(S_j) \to Child(S_k).
\]

This structure makes it possible to reconstruct latency contributions and service 
interactions. The duration of a span is computed as:

\[
Duration(S_k) = EndTime(S_k) - StartTime(S_k).
\]

By aggregating such relationships across many traces, the system can infer how 
services depend on one another during normal and abnormal execution.

\section{Dependency Graph and Causal Graph}
\label{sec:rca_graphs}

\subsection{Service Dependency Graph Construction}
\label{subsec:rca_dependency_graph}

The \emph{service dependency graph} models the structural communication topology 
of the distributed system, defined as:

\[
G = (\mathcal{S}, \mathcal{E}, W),
\]

where $\mathcal{S}$ is the set of services, $\mathcal{E}$ the set of directed 
inter-service calls, and $W$ a weighting function over edges.

A directed edge $(S_i \to S_j) \in \mathcal{E}$ is created when tracing data 
shows that service $S_i$ invokes service $S_j$. To reflect the operational 
strength of this dependency, edge weights are updated dynamically over time. Let 
$c_{ij}^{(t)}$ denote the number of newly observed calls from $S_i$ to $S_j$ 
during interval $t$. The edge weight is updated using exponential decay:

\[
w_{ij}^{(t+1)} = \lambda \, w_{ij}^{(t)} + c_{ij}^{(t)},
\]

where $\lambda \in (0,1)$ is a decay factor that reduces the influence of stale 
observations.

In addition, latency statistics can be associated with each edge. If 
$d_{ij}^{(k)}$ denotes the duration of the $k$-th call from $S_i$ to $S_j$, then 
the average latency is:

\[
\bar{d}_{ij} = \frac{1}{N_{ij}} \sum_{k=1}^{N_{ij}} d_{ij}^{(k)},
\]

where $N_{ij}$ is the number of observed calls. The dependency graph therefore 
captures both structural and operational properties of the service topology.

\subsection{Causal Graph}
\label{subsec:rca_causal_graph}

While the dependency graph captures communication structure, the \emph{causal 
graph} captures the plausible propagation of anomalous behavior through that 
structure. In this framework, the causal graph is derived from the dependency 
graph and conditioned on the current anomaly state of services.

Let the anomaly state of service $S_i$ at time $t$ be defined as:

\[
a_i(t) =
\begin{cases}
1, & \text{if } S_i \text{ is anomalous at time } t, \\
0, & \text{otherwise.}
\end{cases}
\]

The set of anomalous services is therefore:

\[
\mathcal{A}(t) = \{S_i \in \mathcal{S} \mid a_i(t)=1\}.
\]

A causal interpretation is then built by reasoning over graph directionality and 
topology: anomalies are assumed to propagate only along observed dependency 
edges, and upstream services are treated as more plausible sources than 
downstream ones.

For a given service $S_i$, its upstream neighborhood within $h$ hops is defined 
as:

\[
\mathcal{U}_h(S_i) = \{S_j \in \mathcal{S} \mid \exists \text{ path } S_j \leadsto S_i,\ \text{length} \le h\},
\]

while its downstream neighborhood is:

\[
\mathcal{D}_h(S_i) = \{S_j \in \mathcal{S} \mid \exists \text{ path } S_i \leadsto S_j,\ \text{length} \le h\}.
\]

These neighborhoods provide the structural basis for distinguishing likely root 
causes from likely propagated anomalies.

\subsection{Root Cause Indicators in Causal Graphs}
\label{subsec:rca_indicators}

Several graph-based patterns indicate plausible root causes.

First, a service is more likely to be a root cause if it has few anomalous 
upstream dependencies. This can be expressed as:

\[
\phi(S_i) = \frac{|\mathcal{U}_h(S_i) \cap \mathcal{A}(t)|}{|\mathcal{U}_h(S_i)|},
\]

with smaller values of $\phi(S_i)$ suggesting that the anomaly is less likely to 
have been inherited from upstream services.

Second, a root cause candidate should explain meaningful downstream impact. One 
possible impact score is:

\[
I(S_i) = |\mathcal{D}_h(S_i)| + \sum_{S_j \in \mathcal{N}^+(S_i)} w_{ij},
\]

where $\mathcal{N}^+(S_i)$ denotes the set of immediate downstream neighbors.

Finally, temporal ordering also provides useful evidence. If anomalies in $S_i$ 
consistently precede those in downstream services $S_j$, then $S_i$ becomes a 
more plausible causal source:

\[
t_a(S_i) < t_a(S_j).
\]

Together, these indicators support a structural and temporal interpretation of 
failure propagation.

\section{Root Cause Scoring}
\label{sec:rca_scoring}

\subsection{PageRank-Based Root Cause Scoring}
\label{subsec:rca_pagerank}

To prioritize likely root causes, the framework adopts a graph-based ranking 
approach inspired by PageRank. This is appropriate because root causes are not 
only directly connected to anomalous services, but may also influence them 
indirectly through multi-hop dependency paths \cite{brin1998anatomy, fu2020root}.

Let $A \in \mathbb{R}^{n \times n}$ denote the weighted adjacency matrix of the 
causal graph, where:

\[
A_{ij} =
\begin{cases}
w_{ij}, & \text{if } (S_i \to S_j) \in \mathcal{E}, \\
0, & \text{otherwise.}
\end{cases}
\]

After row normalization, the transition matrix becomes:

\[
P_{ij} = \frac{A_{ij}}{\sum_{k=1}^{n} A_{ik}}.
\]

The personalized PageRank vector $\mathbf{r}$ is then defined as:

\[
\mathbf{r} = (1-\gamma)\mathbf{v} + \gamma P^\top \mathbf{r},
\]

where $\gamma$ is the damping factor and $\mathbf{v}$ is a personalization vector 
biased toward anomalous or candidate root cause services.

To combine structural influence with operational impact, a final RCA score may be 
computed as:

\[
Score(S_i) = \alpha \, r_i + \beta \, \hat{I}(S_i),
\]

where $r_i$ is the PageRank score of service $S_i$, $\hat{I}(S_i)$ is a 
normalized impact score, and $\alpha + \beta = 1$.

This formulation favors services that are both graph-central and capable of 
explaining broad downstream degradation.

\subsection{Ranking and Presenting Results}
\label{subsec:rca_ranking}

The final RCA output is a ranked list of candidate root causes:

\[
\text{Rank}(S_i) = \text{argsort}_{\text{desc}}(Score(S_i)).
\]

Rather than returning a single deterministic answer, the system presents the 
highest-ranked services together with supporting evidence, such as downstream 
impact and candidate propagation paths. For a root cause candidate $S_i$, the set 
of explanatory paths may be written as:

\[
\Pi_i = \{ Path(S_i, S_j) \mid S_j \in \mathcal{A}(t),\ S_i \leadsto S_j \}.
\]

This ranking-and-explanation view is especially useful in operational settings, 
where ambiguity is common and RCA should assist human investigation rather than 
replace it entirely.

\section{Implementation}
\label{sec:rca_implementation}

\subsection{Framework Components}
\label{subsec:rca_components}

The RCA engine is composed of four principal components.

\textbf{Trace Collector.} An ingestion module subscribes to distributed trace
data produced by the OpenTelemetry exporter running alongside each monitored
microservice. Incoming traces are decomposed into individual spans. For each
parent--child span pair, the collector emits a dependency record
$(S_i \to S_j, \textit{call\_count}, \textit{latency})$ and forwards it to
the graph builder. The collector runs continuously in a background coroutine
and is resilient to trace gaps: if the upstream exporter is temporarily
unavailable, buffered spans are replayed once connectivity resumes (see
Chapter~\ref{ch:reliability}).

\textbf{Causal Graph Builder (\texttt{CausalGraph}).} Maintains a weighted,
directed graph $G=(\mathcal{S},\mathcal{E},W)$ using the NetworkX library.
New dependency records update edge weights through exponential decay as
described in Section~\ref{subsec:rca_dependency_graph}. The builder also
tracks per-edge call counts and rolling average latencies (retaining the most
recent~100 observations). Weak edges whose weight falls below a configurable
minimum $w_{\min}=0.01$ are pruned periodically, keeping the graph faithful to
the current operational topology. Anomalous services are recorded via a
\texttt{mark\_anomaly()} method that annotates the corresponding graph node
with an \texttt{is\_anomalous} flag and a timestamp.

\textbf{PageRank Scorer (\texttt{RootCauseAnalyzer}).} Upon receiving an
anomaly event, the scorer (i) subsets the causal graph to the anomalous node
set $\mathcal{A}(t)$, (ii) classifies each anomalous service as either a
\emph{root cause} or a \emph{propagated anomaly} by inspecting upstream
dependencies, and (iii) computes the impact score $I(S_i)$ defined in
Section~\ref{subsec:rca_indicators}. Services whose anomalous upstream
neighbourhood is empty---i.e., for which no anomalous predecessor exists in the
graph---are labeled as root causes, while the rest are labeled as propagated.
The output is a ranked list of root-cause candidates sorted by impact score.
Benchmark experiments show that this approach achieves a Top-1 accuracy of 42\%
and a Top-3 accuracy of 89\% over 100~fault-injection events in the
Train-Ticket topology, and scales effectively to a 50-service Alibaba cluster
trace topology (Section~\ref{sec:eval_rca}).

\textbf{Persistent Storage.} The graph state is periodically snapshotted to a
JSON structure containing all nodes, edges with weights and call counts, and
the current anomalous-service set. Snapshots are persisted to durable storage
with a configurable interval (default~300\,s), ensuring that the graph
survives coordinator restarts without requiring a full re-ingestion of trace
data.

\subsection{Integration with Anomaly Detection}
\label{subsec:rca_integration}

The RCA engine is tightly integrated with the anomaly detection pipeline
through the following event-driven workflow.

\begin{enumerate}
    \item \textbf{Detection Trigger.} When the LSTM-AE running on an edge node
    produces a reconstruction error $e_t^{(i)}$ that exceeds the current
    adaptive threshold $\tau_t^{(i)}$
    (Chapter~\ref{ch:adaptive_thresholding}), the node emits an anomaly event
    containing the service name, timestamp, anomaly score, and the current
    threshold value.

    \item \textbf{Anomaly Registration.} The coordinator receives the event
    and calls \\ \texttt{causal\_graph.mark\_anomaly(service)}, which annotates
    the corresponding node in the dependency graph. If multiple services fire
    within a configurable time window (default~60\,s), their events are batched
    into a single anomalous set $\mathcal{A}(t)$.

    \item \textbf{RCA Query.} Once the batching window closes, the coordinator
    invokes \texttt{analyzer.analyze($\mathcal{A}(t)$)}. The analyzer inspects
    the causal graph, classifies services as root causes or propagated
    anomalies, and returns a ranked list of candidates together with metadata
    such as impact scores, downstream service counts, and propagation paths.

    \item \textbf{Alert Enrichment.} The RCA results are attached to the
    original anomaly alert. The enriched alert contains: (a) the ranked
    root-cause candidates with scores, (b) the set of propagated services and
    their upstream explanatory paths, and (c) graph statistics (node count,
    edge count, staleness since last trace update). This enriched alert is
    forwarded to the alerting and dashboard subsystem, providing operators with
    actionable information beyond a simple anomaly flag.
\end{enumerate}

This integration ensures that every anomaly alert is accompanied by a causal
explanation, reducing the investigative burden on operators and directly
supporting MTTR reduction.
